\begin{pastexam}[label={px:ntu-112-bime-ds-01}]{臺大生機所 112 學年度（A）第 1 題}
（問答）

(20\%) (a) Please explain what \textbf{2-3 trees} are. (b) Please list the advantages of using
\textbf{red-black trees} to implement the concept of 2-3 trees.
\end{pastexam}

\begin{pastsoln}
\textbf{(a) 2-3 樹}（資結書第~11.1~節）是一種多路搜尋樹，滿足：
\begin{enumerate}
\item 每個內部節點是 \textbf{2-節點}（$1$ 個鍵、$2$ 個子節點）或 \textbf{3-節點}（$2$ 個鍵
$a < b$、$3$ 個子節點）。
\item 搜尋樹的順序：2-節點左子樹的鍵都小於 $a$、右子樹都大於 $a$；3-節點的三棵子樹分別放
$< a$、介於 $a$ 和 $b$ 之間、$> b$ 的鍵。
\item 葉節點放 $1$ 或 $2$ 個鍵，而且\textbf{所有葉節點在同一層}。
\end{enumerate}
所有葉節點同層，樹就一定是平衡的。 高度 $h$ 的 2-3 樹，鍵最少時全是 2-節點，有 $2^h - 1$ 個鍵；
最多時全是 3-節點，有 $3^h - 1$ 個鍵。 所以 $n$ 個鍵時
$\log_3(n + 1) \le h \le \lg(n + 1)$，搜尋、插入、刪除都是 $O(\log n)$。 插入時把鍵放進葉節點；
葉節點變成 $3$ 個鍵就從中間\textbf{分裂}，中間的鍵往上併進父節點，必要時一路分裂到根，根分裂時
整棵樹長高一層。 樹只從根長高，所以葉節點永遠同層。

\textbf{(b) 用紅黑樹實作 2-3 樹的好處。} 對應的方式（資結書第~11.1.6~節）：2-節點是一個黑色節點；
3-節點 $\{a, b\}$ 拆成兩個二元節點，用一條紅色的連結連起來（例如黑色的 $b$ 帶紅色的左子節點
$a$）。 這樣 2-3 樹就變成一棵普通的二元搜尋樹，每個節點多一個顏色位元。 好處：
\begin{itemize}
\item \textbf{節點只有一種}：直接實作 2-3 樹要處理 2-節點、3-節點兩種節點，插入時還有暫時的
4-節點，節點之間要互相轉換，程式繁瑣，也常浪費空間（每個節點都要預留兩個鍵、三個指標）。
紅黑樹每個節點都是一樣的二元節點，只多一個位元記顏色。
\item \textbf{沿用二元搜尋樹的程式}：搜尋、最小／最大、前驅／後繼、中序走訪、範圍查詢、第 $k$
小等操作都不看顏色，和普通 BST 的程式完全一樣，不用另外寫。
\item \textbf{仍然保證平衡}：2-3 樹的每一層在紅黑樹裡最多佔兩層（黑加紅），所以高度不超過
$2\lg(n + 1)$，所有操作最壞 $O(\log n)$。
\item \textbf{調整是局部的}：2-3 樹的分裂、合併，在紅黑樹裡變成幾次旋轉和改顏色，每次 $O(1)$。
插入最多旋轉 $2$ 次，刪除最多 $3$ 次，比 AVL 樹的調整少（資結書第~10.5~節）。
\end{itemize}
\end{pastsoln}
